\documentclass[10pt,a4paper]{amsart}
\usepackage[margin=19mm]{geometry}
\usepackage{amsmath,amssymb,mathtools}
\usepackage[scr=rsfs,cal=euler]{mathalfa}
\usepackage{xfrac}
\usepackage[unicode,hidelinks,bookmarks=false]{hyperref}
\newcommand{\ie}{i.e.\ }
\newcommand{\cf}{cf.\ }
\newcommand{\resp}{respectively\ }
\newcommand{\Subsection}{\subsection}
\providecommand{\nobreakdash}{\nobreak\hbox{-}}
\setlength{\emergencystretch}{2em}
\multlinegap0pt
\usepackage{bm}
\usepackage{bbm}
\usepackage{dsfont}
\usepackage{xspace}
\usepackage{cancel}
\usepackage{mathtools}
\usepackage{amsthm}
\makeatletter
\@ifundefined{definition}{\newtheorem{definition}{\bf Definition}}{}
\@ifundefined{theorem}{\newtheorem{theorem}{\bf Theorem}}{}
\makeatother
\title[Stereographic Projections on Some Quadric Surfaces]{Stereographic Projections on Some Quadric Surfaces}
\author{W.F.C. Barboza, T.F. Cruz, R.B. Leal}
\date{Source: arXiv:2506.08289v1 (CC BY 4.0)}
\begin{document}
\maketitle

\noindent\emph{Source and licence.} Stereographic Projections on Some Quadric Surfaces. Version arXiv:2506.08289v1 (CC BY 4.0), distributed under the Creative Commons Attribution 4.0 International licence, as recorded in the arXiv licence metadata for this exact version and shown on the abstract page https://arxiv.org/abs/2506.08289v1. Full rights notice: licenses/arxiv-2506.08289-CC-BY-4.0.txt.\par\smallskip

\noindent\textit{Editorial scope.} Counted unit: the Theorem whose source label is \texttt{teorema3} in the article (source0.tex lines 787--792, source label \texttt{teorema3}), delivered together with the article's own complete original proof of it (source0.tex lines 794--857, one \texttt{proof} environment of 64 lines). No mathematics is written, repaired or re-derived: statement and proof are the article's own text line for line. Required context retained uncounted: definicao curvatura (lines 133--138); ellipse\_intersection (lines 627--629); elipsenoplanoxy (lines 686--688). Every definition, display and label of those lines is retained unchanged. Omitted from this package and not redistributed: 1--132, 139--626, 630--685, 689--786, 793--793, 858--1047. Nothing inside a retained excerpt is omitted or altered: every retained body is delivered exactly as the article's lines are written, so no omission receipt of any kind accompanies this package. Citations: the retained excerpts carry no citation command at all, so no bibliography entry of the article is transcribed and no reference is redistributed. Reference adaptation: none was needed and none was made; every reference inside the delivered unit resolves inside it, so no unresolved reference or citation is produced. Printed slips kept literal: no printed word, symbol, hypothesis or number of the counted statement or of its proof was altered. Layout and class: the article's own class is \texttt{amsart} with packages including graphicx, pstricks, pst-plot, hyperref, color, amsfonts, babel, amscd, amsmath, amssymb, float, fancyhdr, geometry; the wrapper is the lane's pinned \texttt{amsart} editorial wrapper at 10pt on A4, which restates the theorem-like environments the retained text needs and adds the article's own macro definitions that the retained text needs (none), with extra wrapper packages (bm, bbm, dsfont, xspace, cancel, mathtools). The delivered numbering is the unit's own. Assets: no retained span contains a figure environment, a graphics-inclusion command, a PDF-page inclusion, a table or a TikZ picture; no image file is copied into this package, the package declares no asset dependency and no image file is redistributed with it. Licence: version arXiv:2506.08289v1 of this article is distributed under the Creative Commons Attribution 4.0 International licence, as recorded in the arXiv licence metadata for this exact version and shown on the abstract page https://arxiv.org/abs/2506.08289v1. A full rights notice is delivered at licenses/arxiv-2506.08289-CC-BY-4.0.txt. Characters: every retained excerpt is pure ASCII, so the delivered engine, XeLaTeX with its default Latin Modern text font, reports no missing character.
\begin{definition}
	Let $\alpha: I \longrightarrow \mathbb{R}^2$ be a plane curve $\alpha(t)=(x(t), y(t)),$ the signed curvature of $\alpha$ at $t$ is given by
	\begin{equation}\label{definicao curvatura}
		k(t) = \dfrac{x'(t) y''(t) - x''(t) y'(t)}{[ (x'(t))^2 + (y'(t))^2 ]^{3/2}}.   
	\end{equation}
\end{definition}

\begin{equation}\label{ellipse_intersection}
	\mathcal{E}_d: \quad \dfrac{x^2}{A^2} + \dfrac{y^2}{B^2} = 1,
\end{equation}

	\begin{equation}\label{elipsenoplanoxy}
		\mathcal{E}_0 :   \dfrac{x^2}{(A_0)^2} + \dfrac{y^2}{(B_0)^2} = 1.
	\end{equation}

\begin{theorem}\label{teorema3}
	Let $k_{\mathcal{E}_d}$ and $k_{\mathcal{E}_0}$ be the curvatures of ellipses \eqref{ellipse_intersection} and \eqref{elipsenoplanoxy}, then
	$$
	k_{\mathcal{E}_0}(t) =  \left(\dfrac{c-d}{c} \right) \cdot k_{\mathcal{E}_d}(t).
	$$
\end{theorem}

\begin{proof}
	Let us begin by computing the curvature of the ellipse \( \mathcal{E}_d \). To do this, note that a parameterization for the planar curve  \eqref{ellipse_intersection} is given by
	\begin{equation}\label{parametriza elipse ed}
		\alpha(t) = (A \cos t, B \sin t), \quad t \in [0,2\pi],
	\end{equation}
	where
	$$
	A = \dfrac{a}{c}\sqrt{c^2 - d^2} \quad \text{and} \quad B = \dfrac{b}{c}\sqrt{c^2 - d^2}.
	$$
	Thus, we have that the first derivative gives us the tangent vector
	$$
	\alpha'(t) = (-A \sin t,\, B \cos t),
	$$
	and the second derivative gives us the acceleration vector
	$$
	\alpha''(t) = (-A \cos t, -B \sin t).
	$$
	From definition \ref{definicao curvatura}, the curvature $k_{\mathcal{E}_d}(t)$ of a planar curve $\mathcal{E}_d$ is given by
	\begin{align*}
		k_{\mathcal{E}_d}(t) & = \dfrac{AB \sin^2(t) + AB \cos^2(t)}{A^2 \sin^2(t) + B^2 \cos^2(t)}  = \dfrac{AB}{A^2 \sin^2(t) + B^2 \cos^2(t)}.
	\end{align*}
	Substituting the values of $A$ and $B$, we obtain
	\begin{align*}
		k_{\mathcal{E}_d}(t) &= \frac{\dfrac{ab(c^2 - d^2)}{c^2}}{\left( \dfrac{(c^2 - d^2)}{c^2} \right)^{3/2} (a^2 \sin^2 t + b^2 \cos^2 t)^{3/2}} \\
		&= \frac{ab}{\dfrac{(c^2 - d^2)^{1/2}}{c}} \cdot \frac{1}{(a^2 \sin^2 t + b^2 \cos^2 t)^{3/2}} \\
		&= \frac{abc}{\sqrt{c^2 - d^2}} \cdot \frac{1}{(a^2 \sin^2 t + b^2 \cos^2 t)^{3/2}}.
	\end{align*}
	We conclude that the curvature of \eqref{ellipse_intersection} at a given instant \( t \) is given by
	\begin{equation}\label{curvatura1}
		k_{\mathcal{E}_d}(t) = \frac{abc}{\sqrt{c^2 - d^2}} \cdot \frac{1}{(a^2 \sin^2 t + b^2 \cos^2 t)^{3/2}}.
	\end{equation}
	Analogously to what we did for the ellipse $\mathcal{E}_d$, we can parametrize the ellipse $\mathcal{E}_0$ given in \eqref{elipsenoplanoxy} by
	\begin{equation}\label{parametriza elipse e0}
		\beta(t) = (A_0 \cos t, B_0 \sin t), \quad t \in [0,2\pi],
	\end{equation}
	where,
	$$
	A_0 = \dfrac{a \sqrt{c^2 - d^2}}{c - d} \quad \text{and} \quad B_0 = \dfrac{b \sqrt{c^2 - d^2}}{c - d}.
	$$
	Defining $k_{\mathcal{E}_0}(t)$ as the curvature of this ellipse $\mathcal{E}_0$ at time $t$, we have
	$$
	k_{\mathcal{E}_0}(t) = \frac{A_0 B_0}{\left( (A_0)^2 \sin^2 t + (B_0)^2 \cos^2 t \right)^{3/2}},
	$$
	substituting the values of $A_0$ and $B_0$, we obtain
	\begin{align*}
		k_{\mathcal{E}_0}(t) &= \frac{\dfrac{ab(c^2 - d^2)}{(c - d)^2}}{\left( \dfrac{a^2(c^2 - d^2)}{(c - d)^2}\sin^2 t + \dfrac{b^2(c^2 - d^2)}{(c - d)^2}\cos^2 t \right)^{3/2}}\\
		&=  \frac{\dfrac{ab(c^2 - d^2)}{(c - d)^2}}{\left( \dfrac{c^2 - d^2}{(c - d)^2} (a^2 \sin^2 t + b^2 \cos^2 t) \right)^{3/2}},
	\end{align*}
	thus, we have
	\begin{align*}
		k_{\mathcal{E}_0}(t) &= \frac{ab(c^2 - d^2)}{(c - d)^2} \cdot \frac{1}{\left( \dfrac{c^2 - d^2}{(c - d)^2} \right)^{3/2}} \cdot \frac{1}{(a^2 \sin^2 t + b^2 \cos^2 t)^{3/2}}\\
		&= \frac{ab(c^2 - d^2)}{(c - d)^2} \cdot \frac{(c - d)^3}{(c^2 - d^2)^{3/2}} \cdot \frac{1}{(a^2 \sin^2 t + b^2 \cos^2 t)^{3/2}}.
	\end{align*}
	Therefore, from \eqref{curvatura1} we get 
	\begin{align*}
		k_{\mathcal{E}_0}(t) &= \frac{ab(c - d)}{\sqrt{c^2 - d^2}} \cdot \frac{1}{(a^2 \sin^2 t + b^2 \cos^2 t)^{3/2}}\\
		& = k_{\mathcal{E}_d}(t) - \left(  \dfrac{abd}{\sqrt{c^2 - d^2} \cdot (a^2 \sin^2 t + b^2 \cos^2 t)^{3/2} } \right).
	\end{align*}
	By simplifying a bit further, we can multiply the expression in parentheses by $\dfrac{c}{c} = 1$, and thus we obtain
	$$
	k_{\mathcal{E}_0}(t) = k_{\mathcal{E}_d}(t) - k_{\mathcal{E}_d}(t)\frac{d}{c} = \left(\frac{c-d}{c}\right) \cdot k_{\mathcal{E}_d}(t).
	$$
	
\end{proof}

\end{document}
